\documentclass[11pt]{article}

% This template is designed for XeLaTeX.
\usepackage{geometry}
\geometry{
  a4paper,
  margin=25.4mm
}

\makeatletter
\def\input@path{{../Style/}}
\makeatother

\usepackage{coursework-report-core}

\showlayouttrue
\setlength{\parindent}{20pt}
\setlength{\parskip}{0.5em}

% Optional bibliography file for BibLaTeX/Biber.
\addbibresource{references.bib}

% Set the default level used by \smartsubsection.
% Valid values: section, subsection, subsubsection, paragraph, subparagraph
\setsmartsubsectionlevel{subsection}
% Optional depth offsets:
% \smartsubsection{Subsection}
% \smartsubsection[1]{Subsubsection}
% \smartsubsection[2]{Paragraph}

% Override these if you want a different numbering style for deeper headings.
% Example:
% \renewcommand{\paragraphnumberformat}{\thesubsection.\arabic{paragraph}}
% \renewcommand{\subparagraphnumberformat}{\theparagraph.\arabic{subparagraph}}




\begin{document}

\pagenumbering{Roman}
\tableofcontents
%\printnomenclature
\newpage
\pagenumbering{arabic}

\begin{algorithm}[h!]
  \caption{Face-Based Finite-Volume Solver Loop for a General Shallow-Water Mesh}
  \label{alg:fv_solver_loop}
  \begin{algorithmic}[1]
      \Require Cell states $\mathbf{U}_P^n = [h_P^n, q_{x,P}^n, q_{y,P}^n]^T$, mesh geometry, bathymetry $b_P$, timestep $\Delta t$
      \Ensure Updated cell states $\mathbf{U}_P^{n+1}$ and diagnostic outputs

      \State Initialise residuals for all cells:
      \[
          \mathbf{R}_P \gets \mathbf{0}
          \label{eq:init_residual}
      \]

      \For{each cell $P$}
          \State Recover primitive variables:
          \[
              u_P =
              \frac{q_{x,P}}{h_P},
              \qquad
              v_P =
              \frac{q_{y,P}}{h_P},
              \qquad
              \eta_P =
              h_P + b_P
              \label{eq:primitive_recovery_algorithm}
          \]
          \State Apply wet/dry state logic using $h_{\min}$
      \EndFor

      \For{each face $f$}
          \State Identify owner cell $P_f$ and neighbour/ghost state $N_f$
          \State Reconstruct left and right states at the face:
          \[
              \mathbf{U}_{L,f},
              \qquad
              \mathbf{U}_{R,f}
              \label{eq:face_reconstructed_states}
          \]
          \State Apply hydrostatic reconstruction to obtain:
          \[
              \mathbf{U}_{L,f}^{*},
              \qquad
              \mathbf{U}_{R,f}^{*}
              \label{eq:hydrostatic_reconstructed_states_algorithm}
          \]
          \State Compute the numerical normal flux:
          \[
              \widehat{\mathbf{F}}_{n,f}
              =
              \mathcal{F}
              \left(
                  \mathbf{U}_{L,f}^{*},
                  \mathbf{U}_{R,f}^{*},
                  \mathbf{n}_f
              \right)
              \label{eq:numerical_normal_flux_algorithm}
          \]
          \State Apply positivity or draining-time flux limiting where required

          \State Accumulate the owner-cell residual:
          \[
              \mathbf{R}_{P_f}
              \gets
              \mathbf{R}_{P_f}
              -
              A_f\widehat{\mathbf{F}}_{n,f}
              \label{eq:owner_residual_update_algorithm}
          \]

          \If{$f$ is an internal face}
              \State Accumulate the neighbour-cell residual:
              \[
                  \mathbf{R}_{N_f}
                  \gets
                  \mathbf{R}_{N_f}
                  +
                  A_f\widehat{\mathbf{F}}_{n,f}
                  \label{eq:neighbour_residual_update_algorithm}
              \]
          \EndIf
      \EndFor

      \For{each cell $P$}
          \State Assemble the finite-volume right-hand side:
          \[
              \mathcal{L}(\mathbf{U}_P)
              =
              \frac{\mathbf{R}_P}{|\Omega_P|}
              +
              \mathbf{S}_P
              \label{eq:fv_rhs_algorithm}
          \]
          \State Advance the state using the selected time integrator:
          \[
              \mathbf{U}_P^{n+1}
              =
              \mathcal{T}
              \left(
                  \mathbf{U}_P^n,
                  \mathcal{L}(\mathbf{U}_P),
                  \Delta t
              \right)
              \label{eq:time_update_operator_algorithm}
          \]
          \State Apply post-update wet/dry and positivity corrections
          \State Update cell diagnostics and output quantities
      \EndFor

  \end{algorithmic}
\end{algorithm}

\begin{equation}
  \label{eq:cell_data_container}
  \mathcal{C}_P
  =
  \begin{cases}
      \mathcal{G}_P
      =
      \left\{
          I_P,
          \mathbf{x}_P,
          |\Omega_P|,
          \partial\Omega_P,
          \mathcal{N}(P),
          \tau_P,
          \chi_P
      \right\},
      & \text{cell geometry}
      \\[0.75em]
      \mathcal{U}_P
      =
      \left\{
          h_P,
          q_{x,P},
          q_{y,P}
      \right\},
      & \text{conserved shallow-water state}
      \\[0.75em]
      \mathcal{B}_P
      =
      \left\{
          b_P,
          \eta_P,
          \nabla b_P,
          b_{\min,P},
          b_{\max,P}
      \right\},
      & \text{bathymetry and free-surface data}
      \\[0.75em]
      \mathcal{M}_P
      =
      \left\{
          n_P,
          f_P,
          \nu_{H,P},
          h_{\min,P},
          \kappa_P
      \right\},
      & \text{local source-model parameters}
      \\[0.75em]
      \mathcal{R}_P
      =
      \left\{
          R_{h,P},
          R_{q_x,P},
          R_{q_y,P},
          \dot{h}_P,
          \dot{q}_{x,P},
          \dot{q}_{y,P}
      \right\},
      & \text{residual and right-hand-side data}
      \\[0.75em]
      \mathcal{D}_P
      =
      \left\{
          h_{\max,P},
          \eta_{\max,P},
          |\mathbf{u}|_{\max,P},
          t_{\mathrm{arrival},P},
          t_{\mathrm{wet},P},
          \mathrm{CFL}_P
      \right\},
      & \text{diagnostic and visualisation data}
  \end{cases}
\end{equation}

\begin{equation}
  \label{eq:face_data_container}
  \mathcal{F}_f
  =
  \begin{cases}
      \mathcal{I}_f
      =
      \left\{
          I_f,
          P_f,
          N_f
      \right\},
      & \text{face indexing and connectivity}
      \\[0.75em]
      \mathcal{G}_f
      =
      \left\{
          \mathbf{x}_f,
          A_f,
          \mathbf{n}_f,
          \mathbf{t}_f
      \right\},
      & \text{face geometry}
      \\[0.75em]
      \mathcal{S}_f
      =
      \left\{
          \mathbf{U}_{L,f},
          \mathbf{U}_{R,f},
          \mathbf{U}_{L,f}^{*},
          \mathbf{U}_{R,f}^{*}
      \right\},
      & \text{left, right, and reconstructed states}
      \\[0.75em]
      \mathcal{Q}_f
      =
      \left\{
          u_{n,L},
          u_{n,R},
          c_L,
          c_R,
          S_L,
          S_R
      \right\},
      & \text{wave-speed and normal-flow data}
      \\[0.75em]
      \mathcal{H}_f
      =
      \left\{
          b_{L,f},
          b_{R,f},
          b_f^{*},
          h_{L,f}^{*},
          h_{R,f}^{*}
      \right\},
      & \text{hydrostatic reconstruction data}
      \\[0.75em]
      \mathcal{R}_f
      =
      \left\{
          \widehat{\mathbf{F}}_{n,f},
          A_f\widehat{\mathbf{F}}_{n,f},
          \theta_f
      \right\},
      & \text{normal flux and limited flux contribution}
      \\[0.75em]
      \mathcal{B}_f
      =
      \left\{
          \beta_f,
          \gamma_f,
          \sigma_f
      \right\},
      & \text{boundary, aperture, and mask data}
  \end{cases}
\end{equation}

Here, $I_P$ and $I_f$ denote cell and face indices, respectively; $\mathbf{x}_P$ and $\mathbf{x}_f$ are the cell and face centroids; $|\Omega_P|$ is the cell area; $A_f$ is the face length in two dimensions; $\mathbf{n}_f$ is the outward unit normal from owner cell $P_f$ to neighbour cell $N_f$; and $\mathbf{t}_f$ is the corresponding tangential vector. The variable $\tau_P$ denotes the cell topology, such as triangular, quadrilateral, polygonal, or cut-cell, while $\chi_P$ stores the cell-region classification, such as ocean, land, coastline, harbour, or defence region. The set $\mathcal{N}(P)$ stores the neighbouring cells connected to $P$ through its faces.

The conserved state is stored as $\mathcal{U}_P = \{h_P,q_{x,P},q_{y,P}\}$, where $h_P$ is the water depth and $q_{x,P}=h_Pu_P$, $q_{y,P}=h_Pv_P$ are the depth-integrated momenta. The bathymetric block $\mathcal{B}_P$ stores the bed elevation $b_P$, free-surface elevation $\eta_P=h_P+b_P$, bed-gradient data, and local bed extrema used by wetting, drying, and hydrostatic reconstruction. The source-model block $\mathcal{M}_P$ stores Manning roughness $n_P$, Coriolis parameter $f_P$, horizontal eddy viscosity $\nu_{H,P}$, wet/dry threshold $h_{\min,P}$, and a roughness or material class $\kappa_P$.

The residual block $\mathcal{R}_P$ stores the accumulated finite-volume flux imbalance and the resulting right-hand-side components used by the time integrator. The diagnostic block $\mathcal{D}_P$ stores quantities required for scientific interpretation and numerical-health monitoring, including maximum depth, maximum surface elevation, maximum speed, arrival time, wetting duration, and local CFL number.
% Example smart references:
% \smartref{fig:example}
% \smartref{fig:example_a,fig:example_b}
% \smartref{fig:example, tab:example}

% Example smart subsection:
% \smartsubsection{Configurable subsection heading}
% \smartsubsection[1]{One level deeper}

% Example algorithm:
% \begin{algorithm}
%   \caption{Example algorithm}
%   \label{alg:example}
%   \begin{algorithmic}[1]
%     \RequireWrap{Input data}
%     \EnsureWrap{Processed output}
%     \StateWrap{Long algorithm lines wrap cleanly to the next line.}
%   \end{algorithmic}
% \end{algorithm}

% Example listings:
% \begin{lstlisting}[style=MATLAB, caption={Example MATLAB listing}, label={lst:matlab_example}]
% x = linspace(0, 2*pi, 1000);
% y = sin(x);
% plot(x, y)
% \end{lstlisting}
%
% \begin{lstlisting}[style=Python, caption={Example Python listing}, label={lst:python_example}]
% import numpy as np
% print(np.linspace(0.0, 1.0, 5))
% \end{lstlisting}
%
% \begin{lstlisting}[style=C++, caption={Example C++ listing}, label={lst:cpp_example}]
% #include <iostream>
% int main() {
%     std::cout << "Hello world\n";
%     return 0;
% }
% \end{lstlisting}

% Example landscape table pages:
% \begin{landscapetablepages}
% \begin{longtable}{ll}
% \caption{Wide landscape longtable example.}\\
% \toprule
% A & B \\
% \midrule
% \endfirsthead
% \toprule
% A & B \\
% \midrule
% \endhead
% Value 1 & Value 2 \\
% \bottomrule
% \end{longtable}
% \end{landscapetablepages}

%\printbibliography

\end{document}
